\documentclass[../../main.tex]{subfiles}

\begin{document}

\section{点态收敛和紧致收敛}\label{section:ch07-pointwise-compact-convergence}

空间 $Y^X$ 和 $\mathcal{C}(X,Y)$ 除了一致拓扑之外, 还有其他一些有用的拓扑. 现在我们要讨论其中的三个: 点态收敛拓扑、紧致收敛拓扑以及紧致开拓扑.

\begin{definition}\label{definition:ch07-point-open-topology}
给定集合 $X$ 的一个点 $x$ 以及空间 $Y$ 的一个开集 $U$, 令
\begin{align*}
S(x,U)=\{f\mid f\in Y^X\text{ 和 }f(x)\in U\}.
\end{align*}
所有集合 $S(x,U)$ 构成 $Y^X$ 的拓扑的一个子基, 这个拓扑称为点态收敛拓扑(topology of pointwise convergence)或点开拓扑(point-open topology).
\end{definition}

这个拓扑的一般基元素是子基元素 $S(x,U)$ 的有限交. 因此, 包含函数 $f$ 的一个典型的基元素就是由所有在有限多点处"接近" $f$ 的函数 $g$ 组成. 这样的邻域如\cref{figure:ch07-fig46-1}所示, 它由所有函数图像与所给三个垂直区间都相交的函数构成.

\par
\begin{figure}[htbp]
\centering
\includegraphics[width=0.6\textwidth]{fig46-1.png}
\caption{}\label{figure:ch07-fig46-1}
\end{figure}

$Y^X$ 上的点态收敛拓扑并非新鲜事物, 它就是我们早已研究过的积拓扑. 如果我们用 $J$ 代替 $X$, 将 $J$ 的一般元素记为 $\alpha$, 这样看起来就更加熟悉了. 这时, 使得 $x(\alpha)\in U$ 的所有函数 $x\colon J\to Y$ 构成的集合 $S(\alpha,U)$ 恰好就是 $Y^J$ 的子集 $\pi_\alpha^{-1}(U)$, 而它也正好就是积拓扑的标准基元素.

称其为点态收敛拓扑的理由缘于以下定理:

\begin{theorem}\label{theorem:ch07-pointwise-convergence}
在点态收敛拓扑下, 函数序列 $f_n$ 收敛于函数 $f$, 当且仅当对于每一个点 $x\in X$, $Y$ 中点的序列 $f_n(x)$ 收敛于点 $f(x)$.
\end{theorem}

\begin{proof}
这是\cref{lemma:ch07-product-sequence}所证明的积拓扑中的一个一般性的事实, 在这里只是使用函数空间记号重新陈述一遍.
\end{proof}

\begin{example}\label{example:ch07-xn-not-closed}
考虑空间 $\mathbb{R}^I$, 其中 $I=[0,1]$, 定义为 $f_n(x)=x^n$ 的连续函数序列 $(f_n)$ 在点态收敛拓扑下收敛于函数 $f$, 其中 $f$ 的定义为
\begin{align*}
f(x)=
\begin{cases}
0& \text{对于 }0\leqslant x<1,\\
1& \text{对于 }x=1.
\end{cases}
\end{align*}
这个例子说明, 在点态收敛拓扑下, 连续函数的子空间 $\mathcal{C}(I,\mathbb{R})$ 不是 $\mathbb{R}^I$ 中的闭集.
\end{example}

我们知道, 若一个连续函数序列 $(f_n)$ 在一致拓扑下收敛, 则极限必定是连续的. 然而上面这个例子说明, 一个序列仅在点态收敛拓扑下收敛, 却未必有连续的极限. 人们可能要问, 是否存在某一个拓扑介于这两个拓扑之间, 仍能保证连续函数的收敛序列有连续的极限呢? 答案是肯定的. 只要对空间 $X$ 加一点限制, 而且这个限制还相当宽泛, 即要求 $X$ 是紧致生成的. 如果在以下定义的紧致收敛拓扑下, $(f_n)$ 收敛于 $f$, 就足以保证 $f$ 是连续的了.

\begin{definition}\label{definition:ch07-compact-convergence}
设 $(Y,d)$ 是一个度量空间, $X$ 是一个拓扑空间. 给定 $Y^X$ 的一个元素 $f$, $X$ 的一个紧致子空间 $C$ 以及一个数 $\varepsilon>0$, 令 $B_C(f,\varepsilon)$ 表示 $Y^X$ 中所有满足下式的元素 $g$ 构成的集合:
\begin{align*}
\sup\{d(f(x),g(x))\mid x\in C\}<\varepsilon.
\end{align*}
这些集合 $B_C(f,\varepsilon)$ 组成了 $Y^X$ 的一个拓扑基. 我们称这个拓扑为紧致收敛拓扑(topology of compact convergence) (有时也称它为"紧致集合上的一致收敛拓扑").
\end{definition}

易见这些集合 $B_C(f,\varepsilon)$ 满足作为基的条件. 最关键的一步是注意, 若 $g\in B_C(f,\varepsilon)$, 则对
\begin{align*}
\delta=\varepsilon-\sup\{d(f(x),g(x))\mid x\in C\},
\end{align*}
有 $B_C(g,\delta)\subset B_C(f,\varepsilon)$.

这个拓扑与前面的拓扑不同, 在这个拓扑中包含 $f$ 的基元素由具有下述特点的函数组成: 它不是仅在有限多个点上"接近" $f$, 而是在某一个紧致集的所有点上"接近" $f$.

采用这个术语的合理性来自以下定理. 这个定理的证明是直接的.

\begin{theorem}\label{theorem:ch07-compact-convergence-char}
函数序列 $f_n\colon X\to Y$ 关于紧致收敛拓扑收敛于函数 $f$, 当且仅当对 $X$ 的每一个紧致子空间 $C$, 序列 $f_n\mid C$ 一致收敛于 $f\mid C$.
\end{theorem}

\begin{definition}\label{definition:ch07-compactly-generated}
空间 $X$ 称为紧致生成的(compactly generated), 如果 $X$ 的一个子集 $A$ 满足条件: 对于 $X$ 的每一个紧致子空间 $C$, $A\cap C$ 是 $C$ 中一个开集, 则 $A$ 是 $X$ 中的一个开集.
\end{definition}

这个条件等价于: 如果对于每一个紧致子集 $C$, $B\cap C$ 是 $C$ 中闭集, 那么集合 $B$ 是 $X$ 中闭集. 这是对空间的一种相当宽泛的限制, 因为许多熟知的空间都是紧致生成的. 例如以下引理.

\begin{lemma}\label{lemma:ch07-compactly-generated}
若 $X$ 是局部紧致的, 或者 $X$ 满足第一可数性公理, 则 $X$ 是紧致生成的.
\end{lemma}

\begin{proof}
设 $X$ 是局部紧致的, 对于 $X$ 的每一个紧致子空间 $C$, $A\cap C$ 是 $C$ 的一个开集, 我们证明 $A$ 是 $X$ 的一个开集. 给定 $x\in A$, 选取 $x$ 的一个邻域 $U$, 使它包含于 $X$ 的某一个紧致子空间 $C$ 中. 根据假设, 由于 $A\cap C$ 是 $C$ 的一个开集, 所以 $A\cap U$ 是 $U$ 的一个开集, 从而也是 $X$ 的开集. 于是 $A\cap U$ 就是 $x$ 的一个包含在 $A$ 中的邻域, 所以 $A$ 是 $X$ 的一个开集.

设 $X$ 满足第一可数性公理, 如果对于 $X$ 的每一个紧致子集 $C$, $B\cap C$ 是 $C$ 的一个闭集, 我们证明 $B$ 也是 $X$ 的一个闭集. 设 $x$ 是 $\overline{B}$ 的一个点, 下面证明 $x\in B$. 因为 $X$ 在点 $x$ 处有可数基, 所以存在 $B$ 中点的一个序列 $(x_n)$ 收敛于 $x$. 子空间
\begin{align*}
C=\{x\}\cup\{x_n\mid n\in\mathbb{Z}_+\}
\end{align*}
是紧致的. 因此根据假定 $B\cap C$ 是 $C$ 的一个闭集. 又因为对于每一个 $n$, $B\cap C$ 包含 $x_n$, 所以它也必包含 $x$, 从而 $x\in B$.
\end{proof}

关于紧致生成空间的一个关键性的结论是:

\begin{lemma}\label{lemma:ch07-compactly-generated-continuity}
若 $X$ 是紧致生成的, 则函数 $f\colon X\to Y$ 只要对于 $X$ 的每一个紧致子空间 $C$, $f\mid C$ 是连续的, $f$ 便是连续的.
\end{lemma}

\begin{proof}
设 $V$ 是 $Y$ 的一个开子集, 我们证明 $f^{-1}(V)$ 是 $X$ 中的开子集. 给定 $X$ 的任意一个子空间 $C$,
\begin{align*}
f^{-1}(V)\cap C=(f\mid C)^{-1}(V).
\end{align*}
如果 $C$ 是紧致的, 由 $f\mid C$ 是连续的, 可见这个集合是 $C$ 的一个开集. 因为 $X$ 是紧致生成的, 所以 $f^{-1}(V)$ 是 $X$ 的一个开集.
\end{proof}

\begin{theorem}\label{theorem:ch07-C-closed-compact-convergence}
设 $X$ 是一个紧致生成的空间, $(Y,d)$ 是一个度量空间, 则 $\mathcal{C}(X,Y)$ 关于紧致收敛拓扑是 $Y^X$ 的一个闭集.
\end{theorem}

\begin{proof}
设 $f\in Y^X$ 是 $\mathcal{C}(X,Y)$ 的一个极限点, 我们证明 $f$ 连续. 只需证明, 对于 $X$ 的每一个紧致子集 $C$, $f\mid C$ 是连续的. 对于每一个 $n$, 考虑 $f$ 的邻域 $B_C(f,1/n)$, 因为它和 $\mathcal{C}(X,Y)$ 相交, 所以可以选出一个函数 $f_n\in\mathcal{C}(X,Y)$ 包含在这一邻域中. 函数序列 $f_n\mid C\colon C\to Y$ 一致收敛于函数 $f\mid C$, 因而根据一致极限定理, $f\mid C$ 连续.
\end{proof}

\begin{corollary}\label{corollary:ch07-limit-continuous}
设 $X$ 是一个紧致生成的空间, $(Y,d)$ 是一个度量空间. 如果一个连续函数序列 $f_n\in\mathcal{C}(X,Y)$ 在紧致收敛拓扑下收敛于函数 $f$, 那么 $f$ 是连续的.
\end{corollary}

当 $Y$ 是一个度量空间时, 我们在函数空间 $Y^X$ 上有三个拓扑. 这三者之间的关系可由下述定理来予以说明, 它的证明是直接的.

\begin{theorem}\label{theorem:ch07-topology-inclusion}
设 $X$ 是一个空间, $(Y,d)$ 是一个度量空间. 对于函数空间 $Y^X$ 上的三种拓扑之间有以下包含关系:
\begin{align*}
(\text{一致拓扑})\supset(\text{紧致收敛拓扑})\supset(\text{点态收敛拓扑}).
\end{align*}
如果 $X$ 是紧致的, 则前两个拓扑是一致的. 如果 $X$ 是离散的, 则后两个拓扑是一致的.
\end{theorem}

一致拓扑与紧致收敛拓扑的定义都涉及到空间 $Y$ 上的度量 $d$, 但点态收敛拓扑却不需要. 事实上, 它的定义是针对任意空间 $Y$ 的. 这样自然就会产生疑问: 是否其他两个拓扑中的哪一个可以扩展为对任意拓扑空间 $Y$ 有定义呢? 对于从 $X$ 映入 $Y$ 的所有函数构成的空间 $Y^X$, 这个问题还没有满意的答案. 但是, 对于连续函数空间 $\mathcal{C}(X,Y)$, 却可以得到一些结果. 我们在 $\mathcal{C}(X,Y)$ 上定义一种拓扑, 即紧致开拓扑. 可见若 $Y$ 是一个度量空间, 则这个拓扑就与紧致收敛拓扑一致. 如我们将要看到的, 紧致开拓扑本身十分重要.

\begin{definition}\label{definition:ch07-compact-open}
设 $X$ 和 $Y$ 是两个拓扑空间. 如果 $C$ 是 $X$ 的一个紧致子空间, $U$ 是 $Y$ 的一个开子集, 定义
\begin{align*}
S(C,U)=\{f\mid f\in\mathcal{C}(X,Y)\text{ 和 }f(C)\subset U\}.
\end{align*}
集合 $S(C,U)$ 形成 $\mathcal{C}(X,Y)$ 的某一个拓扑的子基, 这个拓扑就称为紧致开拓扑(compact-open topology).
\end{definition}

根据定义易见, 紧致开拓扑一般比点态收敛拓扑更细, 是可以定义在整个空间 $Y^X$ 上的. 但我们关心的只是子空间 $\mathcal{C}(X,Y)$, 所以只考虑这个空间.

\begin{theorem}\label{theorem:ch07-compact-open-equals-compact-convergence}
设 $X$ 是一个空间, $(Y,d)$ 是一个度量空间. 对于空间 $\mathcal{C}(X,Y)$ 来说, 紧致开拓扑与紧致收敛拓扑一致.
\end{theorem}

\begin{proof}
如果 $A$ 是 $Y$ 的一个子集, $\varepsilon>0$, 我们用 $U(A,\varepsilon)$ 表示 $A$ 的 $\varepsilon$-邻域. 如果 $A$ 是紧致的, $V$ 是包含 $A$ 的一个开集, 于是存在 $\varepsilon>0$, 使得 $U(A,\varepsilon)\subset V$. 事实上, 所求的 $\varepsilon$ 就是函数 $d(a,X-V)$ 的极小值.

首先, 我们证明紧致收敛拓扑比紧致开拓扑细. 设 $S(C,U)$ 是紧致开拓扑的一个子基元素, 而 $f$ 是 $S(C,U)$ 中的一个元素, 由于 $f$ 连续, 所以 $f(C)$ 是开集 $U$ 的一个紧致子集. 因此, 我们可以选取 $\varepsilon$, 使得 $f(C)$ 的 $\varepsilon$-邻域包含在 $U$ 中, 于是
\begin{align*}
B_C(f,\varepsilon)\subset S(C,U).
\end{align*}

我们证明紧致开拓扑比紧致收敛拓扑细. 设 $f\in\mathcal{C}(X,Y)$. 对于紧致收敛拓扑包含 $f$ 的一个开集, 则它包含一个形如 $B_C(f,\varepsilon)$ 的基元素, 下面我们找出紧致开拓扑的一个包含 $f$ 的基元素, 使它包含于 $B_C(f,\varepsilon)$ 中.

$X$ 的每一点 $x$ 都有一个邻域 $V_x$, 使得 $f(\overline{V}_x)$ 包含在 $Y$ 的一个直径小于 $\varepsilon$ 的开集 $U_x$ 中. [例如, 选取 $V_x$, 使 $f(V_x)$ 包含在 $f(x)$ 的 $\varepsilon/4$-邻域中, 则 $f(\overline{V}_x)$ 包含在 $f(x)$ 的 $\varepsilon/3$-邻域中, 这个邻域的直径最多是 $2\varepsilon/3$.] 用有限多个这种集合 $V_x$, 譬如 $x=x_1,\cdots,x_n$ 来覆盖 $C$. 令 $C_x=\overline{V}_x\cap C$, 于是 $C_x$ 为紧致子集, 并且基元素
\begin{align*}
S(C_{x_1},U_{x_1})\cap\cdots\cap S(C_{x_n},U_{x_n})
\end{align*}
包含 $f$, 并且被 $B_C(f,\varepsilon)$ 所包含.
\end{proof}

\begin{corollary}\label{corollary:ch07-compact-open-metric-independent}
设 $Y$ 是一个度量空间. $\mathcal{C}(X,Y)$ 的紧致收敛拓扑与 $Y$ 的度量无关. 因此, 若 $X$ 是紧致的, 则 $\mathcal{C}(X,Y)$ 的一致拓扑与 $Y$ 的度量无关.
\end{corollary}

紧致开拓扑的定义不涉及度量, 这恰是它的一个有用的特征. 同时它还满足"联合连续性"的要求. 粗略地说就是, $f(x)$ 不仅对单变量 $x$ 是连续的, 而且在变量 $x$ 和 $f$ 同时改变时, 仍然是连续的. 确切地说来, 便是下面的定理.

\begin{theorem}\label{theorem:ch07-evaluation-continuous}
设 $X$ 是一个局部紧致的 Hausdorff 空间, $\mathcal{C}(X,Y)$ 有紧致开拓扑. 映射
\begin{align*}
e\colon X\times\mathcal{C}(X,Y)\longrightarrow Y
\end{align*}
定义为
\begin{align*}
e(x,f)=f(x),
\end{align*}
它是连续的.
\end{theorem}

这个映射 $e$ 称为赋值映射(evaluation map).

\begin{proof}
给定 $X\times\mathcal{C}(X,Y)$ 的一个点 $(x,f)$, 设 $V$ 是 $Y$ 中包含像点 $e(x,f)=f(x)$ 的一个开集, 我们希望找到一个包含 $(x,f)$ 的开集, 使得 $e$ 将其映射到 $V$ 中. 首先, 应用 $f$ 的连续性以及 $X$ 是局部紧致的 Hausdorff 空间这一事实, 可以选取一个包含 $x$ 的开集 $U$, 使得 $U$ 有紧致闭包 $\overline{U}$, 并且 $f$ 将 $\overline{U}$ 映射到 $V$ 中. 考虑 $X\times\mathcal{C}(X,Y)$ 中的开集 $U\times S(\overline{U},V)$. 它是包含 $(x,f)$ 的一个开集, 并且若 $(x',f')$ 属于这个集合, 则 $e(x',f')=f'(x')$ 属于集合 $V$.
\end{proof}

这个定理的一个推论如下, 它在代数拓扑中是很有用的.

\begin{definition}\label{definition:ch07-induced-map}
如果给定了一个函数 $f\colon X\times Z\to Y$, 我们便有一个相应的函数 $F\colon Z\to\mathcal{C}(X,Y)$, 其定义为
\begin{align*}
(F(z))(x)=f(x,z).
\end{align*}
反之, 如果给定 $F\colon Z\to\mathcal{C}(X,Y)$, 那么这个等式又定义了与之相应的函数 $f\colon X\times Z\to Y$. 我们称从 $Z$ 到 $\mathcal{C}(X,Y)$ 中的映射 $F$ 是由 $f$ 诱导的(induced by $f$).
\end{definition}

\begin{theorem}\label{theorem:ch07-induced-map}
设 $X$ 和 $Y$ 是两个空间, 赋予 $\mathcal{C}(X,Y)$ 紧致开拓扑. 若 $f\colon X\times Z\to Y$ 是连续的, 则由 $f$ 诱导的映射 $F\colon Z\to\mathcal{C}(X,Y)$ 也是连续的. 进而, 若 $X$ 是局部紧致的 Hausdorff 空间, 则其逆命题也成立.
\end{theorem}

\begin{proof}
我们先假设 $F$ 是连续的, $X$ 是一个局部紧致的 Hausdorff 空间, 那么 $f$ 是连续的. 因为 $f$ 等于复合
\begin{align*}
X\times Z\xrightarrow{i_X\times F}X\times\mathcal{C}(X,Y)\xrightarrow{e}Y,
\end{align*}
其中 $i_X$ 是 $X$ 上的恒等映射.

现在设 $f$ 是连续的, 为了证明 $F$ 的连续性, 我们取 $Z$ 的一个点 $z_0$ 和 $\mathcal{C}(X,Y)$ 中的包含 $F(z_0)$ 的一个子基元素 $S(C,U)$. 下面寻找 $z_0$ 的一个邻域 $W$, 使得 $F$ 可以将 $W$ 映射到 $S(C,U)$ 中. 这样就完成了证明.

所谓 $F(z_0)$ 属于 $S(C,U)$, 即是说, 对任意 $x\in C$, 有 $(F(z_0))(x)=f(x,z_0)$ 属于 $U$. 也就是说 $f(C\times z_0)\subset U$. $f$ 的连续性蕴涵着 $f^{-1}(U)$ 是 $X\times Z$ 中包含 $C\times z_0$ 的一个开集, 因此
\begin{align*}
f^{-1}(U)\cap(C\times Z)
\end{align*}
是 $C\times Z$ 中包含 $C\times z_0$ 的开集. 根据第 26 节中的管状引理可见, 在 $Z$ 中存在 $z_0$ 的一个邻域 $W$, 使得 $C\times W$ 包含于 $f^{-1}(U)$, 如\cref{figure:ch07-fig46-2}所示. 于是对于 $z\in W$ 和 $x\in C$, 有 $f(x,z)\in U$. 从而 $F(W)\subset S(C,U)$.
\end{proof}

\par
\begin{figure}[htbp]
\centering
\includegraphics[width=0.75\textwidth]{fig46-2.png}
\caption{}\label{figure:ch07-fig46-2}
\end{figure}

我们来简略地讨论一下紧致开拓扑与同伦这个概念之间的关系, 其中同伦将在代数拓扑中涉及.

如果 $f$ 和 $g$ 是从 $X$ 到 $Y$ 的两个连续映射, 我们说 $f$ 和 $g$ 是同伦的, 如果存在一个连续映射
\begin{align*}
h\colon X\times[0,1]\longrightarrow Y,
\end{align*}
使得对于所有的 $x\in X$ 有 $h(x,0)=f(x)$ 和 $h(x,1)=g(x)$. 映射 $h$ 称为 $f$ 和 $g$ 之间的一个同伦.

粗略地说, 一个同伦就是从 $X$ 到 $Y$ 的映射构成的一个"连续单参数族". 更准确些讲, 一个同伦 $h$ 给出一个映射
\begin{align*}
H\colon[0,1]\longrightarrow\mathcal{C}(X,Y),
\end{align*}
对于每一个参数 $t\in[0,1]$, 它对应着从 $X$ 到 $Y$ 的一个连续映射. 假设 $X$ 是局部紧致的 Hausdorff 空间, 我们知道, $h$ 连续当且仅当 $H$ 连续. 这就意味着连续映射 $f$ 和 $g$ 之间的一个同伦 $h$, 恰好对应着函数空间 $\mathcal{C}(X,Y)$ 中连接 $\mathcal{C}(X,Y)$ 的点 $f$ 和 $g$ 的一条道路.

我们将在本书的第二部分更详细地研究同伦的概念.

\subsection*{习题}

\begin{enumerate}
\item 证明: 所有集合 $B_C(f,\varepsilon)$ 形成 $Y^X$ 的某一个拓扑的基.
\item 证明\cref{theorem:ch07-topology-inclusion}.
\item 证明有界函数 $f\colon\mathbb{R}\to\mathbb{R}$ 构成的集合 $\mathcal{B}(\mathbb{R},\mathbb{R})$ 关于一致拓扑是 $\mathbb{R}^{\mathbb{R}}$ 的闭集, 但在紧致收敛拓扑下却不是.
\item 考虑定义为
\begin{align*}
f_n(x)=x/n
\end{align*}
的连续函数序列 $f_n\colon\mathbb{R}\to\mathbb{R}$ 在\cref{theorem:ch07-topology-inclusion}中提到的三个拓扑中, 关于哪个拓扑是收敛的? 对于第 21 节的习题 9 中所给定的序列, 情况又如何呢?
\item 考虑定义为
\begin{align*}
f_n(x)=\sum_{k=1}^n kx^k
\end{align*}
的函数序列 $f_n\colon(-1,1)\to\mathbb{R}$.
\begin{enumerate}[(a)]
\item 证明 $(f_n)$ 关于紧致收敛拓扑是收敛的, 并且极限函数是连续的. (这是关于幂级数的一个一般性结论.)
\item 证明 $(f_n)$ 关于一致拓扑不收敛.
\end{enumerate}
\item 证明: 关于紧致开拓扑, 如果 $Y$ 是 Hausdorff 的, 则 $\mathcal{C}(X,Y)$ 也是 Hausdorff 的. 如果 $Y$ 是正则的, 则 $\mathcal{C}(X,Y)$ 也是正则的. [提示: 如果 $\overline{U}\subset V$, 则 $\overline{S(C,U)}\subset S(C,V)$.]
\item 证明: 如果 $Y$ 是局部紧致的 Hausdorff 空间, 只要全都使用紧致开拓扑, 那么复合映射
\begin{align*}
\mathcal{C}(X,Y)\times\mathcal{C}(Y,Z)\longrightarrow\mathcal{C}(X,Z)
\end{align*}
是连续的. [提示: 如果 $g\circ f\in S(C,U)$, 找出 $V$, 使得 $f(C)\subset V$ 和 $g(\overline{V})\subset U$ 成立.]
\item 设 $\mathcal{C}'(X,Y)$ 表示关于某拓扑 $\mathcal{T}$ 的集合 $\mathcal{C}(X,Y)$. 证明: 如果赋值映射
\begin{align*}
e\colon X\times\mathcal{C}'(X,Y)\longrightarrow Y
\end{align*}
连续, 则 $\mathcal{T}$ 包含紧致开拓扑. [提示: 诱导映射 $E\colon\mathcal{C}'(X,Y)\to\mathcal{C}(X,Y)$ 是连续的.]
\item 这是\cref{theorem:ch07-induced-map}在商映射上的一个 (意外的) 应用. (参照第 29 节的习题 11.)
\begin{quote}
\textbf{定理.} 如果 $p\colon A\to B$ 是一个商映射, 并且 $X$ 是局部紧致的 Hausdorff 空间, 那么 $i_X\times p\colon X\times A\to X\times B$ 是一个商映射.
\end{quote}
证明:
\begin{enumerate}[(a)]
\item 设 $Y$ 是由 $i_X\times p$ 诱导的商空间, $q\colon X\times A\to Y$ 是一个商映射. 证明存在一个连续的一一映射 $f\colon Y\to X\times B$, 使得 $f\circ q=i_X\times p$.
\item 设 $g=f^{-1}$, $G\colon B\to\mathcal{C}(X,Y)$ 和 $Q\colon A\to\mathcal{C}(X,Y)$ 分别是由 $g$ 和 $q$ 诱导的映射. 证明 $Q=G\circ p$.
\item 证明 $Q$ 是连续的. 推证 $G$ 是连续的, 因此 $g$ 是连续的.
\end{enumerate}
\item 一个空间称为局部紧致的, 如果存在它的一个开覆盖, 这个开覆盖中的每一个开集都包含于 $X$ 的某一个紧致子集中. 一个空间称为 $\sigma$-紧致的 ($\sigma$-compact), 如果它能被可数多个这样的开集覆盖.
\begin{enumerate}[(a)]
\item 证明: 若 $X$ 是局部紧致的, 并且是第二可数的, 则它是 $\sigma$-紧致的.
\item 设 $(Y,d)$ 是一个度量空间. 证明: 如果 $X$ 是 $\sigma$-紧致的, 那么对 $Y^X$ 上的紧致收敛拓扑, 存在一个度量使得若 $(Y,d)$ 是完备的, 则 $Y^X$ 在此度量下也是完备的. [提示: 设 $A_1,A_2,\cdots$ 是 $X$ 中的可数个紧致子空间构成的族, 并且这些紧致子空间的内部构成了 $X$ 的一个覆盖. 用 $Y_i$ 表示所有从 $A_i$ 到 $Y$ 的映射构成的集合, 并赋予其一致拓扑. 定义 $Y^X$ 和积空间 $Y_1\times Y_2\times\cdots$ 的一个闭子空间之间的一个同胚.]
\end{enumerate}
\item 设 $(Y,d)$ 是一个度量空间, $X$ 是一个空间. 定义 $\mathcal{C}(X,Y)$ 的一个拓扑如下: 给定 $f\in\mathcal{C}(X,Y)$ 和 $X$ 上的一个正的连续函数 $\delta\colon X\to\mathbb{R}_+$, 令
\begin{align*}
B(f,\delta)=\{g\mid d(f(x),g(x))<\delta(x),\text{ 对于所有 }x\in X\}.
\end{align*}
\begin{enumerate}[(a)]
\item 证明这些集合 $B(f,\delta)$ 构成了 $\mathcal{C}(X,Y)$ 的某个拓扑的一个基. 将这个拓扑称为细拓扑(fine topology).
\item 证明细拓扑包含一致拓扑.
\item 若 $X$ 是紧致的, 则细拓扑和一致拓扑是一致的.
\item 证明若 $X$ 是离散的, 则 $\mathcal{C}(X,Y)=Y^X$, 并且细拓扑和箱拓扑一致.
\end{enumerate}
\end{enumerate}

\end{document}